\documentclass[STT-4300-Main.tex]{subfiles}
\title{Chapitre 2 - Modèles linéaires généralisés (MLG)}
\subtitle{STT-4300}
\date{\today}
\begin{document}
\begin{refsection}

\begin{frame}[t]
\titlepage
\end{frame}

\section{Modèles linéaires généralisés (MLG)}

\begin{frame}{Modèles linéaires généralisés (MLG)}
Classe de modèles unifiée par \citeN{Neld:Wedd:1972} (voir \citeNP{McCu:Neld:1989}).

\bigskip

Le bien connu {\color{blue}modèle linéaire} est souvent écrit
$$Y_i = x_i^T \beta + \epsilon_i,$$
avec $\epsilon_i \sim {\mathcal N}(0,\sigma^2)$ et donc $Y_i \sim {\mathcal N}(\mu_i,\sigma^2)$.

Le même modèle peut aussi s'écrire
$$E(Y_i) = \mu_i = x_i^T \beta.$$
\end{frame}

\begin{frame}{}
Les modèles linéaires généralisés (MLG) (\textit{generalized linear models}, GLM) incluent des situations plus générales :
\begin{enumerate}
\item Des variables réponses ayant des distributions autres que la loi
  normale (même discrètes ou catégorielles).
\item La relation entre la réponse et les variables explicatives
  n'a pas besoin d'être de la forme linéaire simple ci-dessus.
\end{enumerate}

La forme générale d'un MLG est :
$${\color{blue}g(E(Y_i)) = g(\mu_i)=\eta_i=x_i^T \beta},$$
où $g$ est appelée la fonction de lien (\textit{link function}). On note $Var(Y_i) = v_i = v(\mu_i)$.

\end{frame}


\begin{frame}{Exemples}

\begin{itemize}
\item régression logistique :\\
 $Y_i \sim  Bernoulli(p_i)$, $E(Y_i) = p_i$
$$\log\Big(\frac{p_i}{1-p_i}\Big) = \mbox{logit}(p_i) = x_i^T \beta.$$
\item  Poisson :\\
$Y_i \sim {\mathcal P}(\lambda_i)$, $E(Y_i) = \lambda_i$
$$\log(\lambda_i) = x_i^T \beta.$$
\end{itemize}

\end{frame}

\begin{frame}{Limitations}

\begin{itemize}
\item la composante linéaire est conservée ( $\leadsto$ MAG, Chapitre~\ref{GAM.section})
\item les distributions sont restreintes à la famille exponentielle (\textit{exponential family}) (voir la discussion dans la
Section~\ref{GLM.estim.section})
\item les réponses doivent être indépendantes ($\leadsto$ GEE, Chapitre~\ref{GEE.section}
et MLGM, Chapitre~\ref{GLMM.section})
\end{itemize}

\end{frame}

\subsection{Famille exponentielle}

\begin{frame}{}


La théorie unifiée des MLG est construite sur la {\color{blue}famille exponentielle} (\textit{exponential family}) de
distributions, tirant avantage de ses bonnes propriétés.

\bigskip

\textbf{Définition} : Considérons $Y$ dont la distribution dépend d'un seul
paramètre $\theta$. Une densité ou fonction de probabilité appartient à
la {\color{blue}famille exponentielle} si elle peut s'écrire sous la forme
$$f(y; \theta,\phi) = \exp\left( A \ \frac{y \theta - b(\theta)}{\phi} +
c(y,\frac{\phi}{A}) \right),$$ 
pour des fonctions spécifiques $b()$ et $c()$ et une constante $A$. Si $\phi$ est
connu, c'est une famille exponentielle avec {\color{blue}paramètre canonique} $\theta$.

\bigskip
Une fonction de lien $g$ de la spécification du MLG est dite canonique si elle est
telle que $\theta(\mu)=\eta$, c'est-à-dire si $g(\cdot) = (b^{\prime})^{-1}(\cdot)$
(étant donné que $\mu=b^{\prime}(\theta)$, voir~\eqref{expect.exponential} à la
page~\pageref{expect.exponential}).   
\end{frame}


\begin{frame}{}

\textbf{Exemple} : La {\color{blue} loi normale} (où $\mu$ est le
paramètre d'intérêt et $\sigma^2$ est considéré comme un paramètre de
nuisance)  appartient à la famille
exponentielle. En effet :
\begin{eqnarray*}
  f(y; \mu) & = & \frac{1}{\sqrt{2\pi}\sigma} \exp\big(-\frac{1}{2
  \sigma^2}(y-\mu)^2\big)\\
& = & \exp \big( -\frac{y^2}{2 \sigma^2} + \frac{y \mu}{\sigma^2} - \frac{\mu^2}{2 \sigma^2} - \frac{1}{2} \log(2 \pi \sigma^2) \big).
\end{eqnarray*}
On identifie $\theta=\mu$, $A=1$, $\phi=\sigma^ 2$, $b(\theta) = \mu^2/2=
\theta^2/2$ et
$$c(y,\phi/A) = -\frac{1}{2} \left( \frac{y^2}{\phi} + \log(2 \pi \phi) \right).$$
Le lien canonique est la fonction identité.
\end{frame}

\begin{frame}{}

\textbf{Exemple} : La {\color{blue}loi de Poisson} appartient à la
famille exponentielle.
\begin{eqnarray*}
  f(y;\lambda) & = & \frac{\lambda^y \exp(-\lambda)}{y!}\\
& = & \exp \big(y \log(\lambda) - \lambda - \log(y!) \big),
\end{eqnarray*}
où $\theta=\log(\lambda)$, $A=1$, $\phi=1$, $b(\theta) = \lambda = \exp(\theta)$
et $c(y,\phi/A)=-\log(y!)$.

\bigskip
Le lien canonique est la fonction log().
\end{frame}


\begin{frame}{}

\begin{table}[tbp]
  \centering
  \begin{tabular}{|l|c|c|c|c|}
   \hline
  Distribution & $A$ & $\phi$ & $b(\theta)$ & $c(y, \phi/A)$\\
  \hline
   ${\mathcal N}(\mu,\sigma^2)$ & $1$ & $\sigma^2$ &  $\theta^2/2$ & $-\frac{1}{2} (y^2/\phi + \log(2 \pi \phi))$ \\
   ${\mathcal P}(\lambda)$ & $1$ & $1$ & $\exp(\theta)$ & $-\log(y!)$ \\
   ${\mathcal B}(m,p)/m$ & & & & \\
    $\Gamma(\mu,\nu)$ & & & &\\
   \hline
  \end{tabular}
  \caption{Membres de la famille exponentielle.}
  \label{exponential.table}
\end{table}
\end{frame}

\begin{frame}{Propriétés de la famille exponentielle}

Développons d'abord des expressions générales pour l'espérance et la variance
de $Y$, utiles plus loin dans ce chapitre.

\bigskip

\textbf{Rappel} : Si $l(\theta, \phi;y) = \log(f(y; \theta,\phi))$, on a :
\begin{equation}
\label{reminder1}
E\left(\frac{\partial l(\theta, \phi;Y)}{\partial \theta}\right)=0,
\end{equation}
et
\begin{equation}
\label{reminder2}
E\left(\frac{\partial^2 l(\theta, \phi;Y)}{\partial \theta^ 2}\right) +
E\left(\left(\frac{\partial l(\theta, \phi;Y)}{\partial \theta}\right)^2\right)=0,
\end{equation}
si l'ordre d'intégration et de dérivation peut être inversé (support compact
de la distribution).
\end{frame}


\begin{frame}{}


Pour les densités de la famille exponentielle, on a
\begin{equation}
\label{deriv1}
\frac{\partial l(\theta, \phi;y)}{\partial \theta} = A \  \frac{y -
b^{\prime}(\theta)}{\phi},  
\end{equation}
et
\begin{equation}
\label{deriv2}
\frac{\partial^2 l(\theta, \phi;y)}{\partial \theta^2} = - A \frac{b^{\prime \prime}(\theta)}{\phi}
\end{equation}

\bigskip
De \eqref{reminder1} et \eqref{deriv1} on obtient
$$0=E\left(\frac{\partial l(\theta, \phi;Y)}{\partial \theta}\right)=A \
\frac{E(Y) - b^{\prime}(\theta)}{\phi}$$
ce qui donne
\bigskip
\begin{equation}
\label{expect.exponential}
{\color{blue}E(Y) = \mu = b^{\prime}(\theta)}
\end{equation}
\end{frame}


\begin{frame}{}

De même, de \eqref{reminder2} et \eqref{deriv2}, et en utilisant
\eqref{expect.exponential}, on obtient

\begin{eqnarray*}
0 & = & E\left(- A \ \frac{b^{\prime \prime}(\theta)}{\phi}\right) +   E\left(
\left( A \ \frac{Y - b^{\prime}(\theta)}{\phi}\right)^2\right)\\
 & = &  - A \ \frac{b^{\prime \prime}(\theta)}{\phi} + A^2  \ E\left(
\left( \frac{Y - \mu}{\phi}\right)^2\right)\\
 & = & - \frac{b^{\prime \prime}(\theta)}{\phi} + A \ \frac{Var(Y)}{\phi^2}\\
 & = & - b^{\prime \prime}(\theta) + A \frac{Var(Y)}{\phi}
\end{eqnarray*}
\bigskip
ce qui donne
\begin{equation}
\label{var.exponential}
{\color{blue}Var(Y) = \frac{b^{\prime \prime}(\theta)\phi}{ A}}
\end{equation}
\end{frame}

\subsection{Composantes d'un MLG}

\begin{frame}{}

{\color{blue}Trois composantes} :

\begin{enumerate}
\item Des variables réponses indépendantes $Y_1, \ldots,Y_n$ supposées partager
  la même distribution de la famille exponentielle ($F_{\theta_i}$).
\item Un ensemble de paramètres $\beta$ et de variables explicatives $x_i^T = (1,
x_{i1},  \ldots, x_{ip})$ pour $i=1,\ldots,n$.
\item Une fonction de lien $g$ monotone et dérivable telle que
$$g(\mu_i) = x_i^T  \beta.$$
$\eta_i = x_i^T \beta$ est appelé le prédicteur linéaire (\textit{linear predictor}).
\end{enumerate}

\end{frame}

\begin{frame}{}

Le problème initial a $n$ (= taille de l'échantillon) paramètres inconnus $\mu_1,
\ldots,\mu_n$ (ou $\theta_1, \ldots, \theta_n$) qui sont réduits à $(p+1)$
paramètres inconnus ($\beta_0, \beta_1,\ldots,\beta_p$) en imposant la structure
d'un modèle MLG.

\bigskip
\bigskip

La transformation via la fonction de lien assure que le
paramètre estimé se situe dans l'espace admissible des valeurs
(par exemple, $(0,1)$ pour Bernoulli/binomiale et $(0, \infty)$ pour Poisson).
\end{frame}

\subsection{Exemples}

\begin{frame}{Loi normale : $Y_i \sim {\mathcal N}(\mu_i,\sigma^2)$}


La loi normale est utilisée pour modéliser des données continues ayant une
distribution symétrique. Exemples : taille ou pression artérielle des individus. La
fonction de lien correspondante est la fonction identité $g(\mu_i)=\mu_i$ (aucune restriction sur les
valeurs possibles de $\mu_i$, un modèle linéaire).

\bigskip
\bigskip

\textbf{Remarque} : Ce modèle est très populaire :
\begin{itemize}
\item beaucoup de phénomènes bien décrits par la loi normale
\item la moyenne ou le total d'un échantillon aléatoire est approximé par la loi normale (théorème central limite)
\item beaucoup de théorie disponible pour ce modèle (résultats explicites, analytiques).
\end{itemize}
\end{frame}

\begin{frame}{Loi binomiale : $Y_i \sim {\mathcal B}(m_i,p_i)$}

Modélise un processus à résultats binaires. Exemples : le nombre de candidats qui
réussissent un test (résultat : réussite ou échec), le nombre de patients atteints d'une maladie
qui sont vivants à un temps spécifié depuis le diagnostic (résultat : vivant ou décédé), le
statut tabagique d'une personne.

\bigskip

Pour les MLG, on préfère considérer la proportion de « succès » $P_i = Y_i/m_i$, pour
laquelle $E(P_i) = \mu_i = p_i$.

\smallskip
Les fonctions de lien couramment utilisées pour la loi binomiale sont (restreignent le
domaine de $\mu_i$ dans (0,1)) :
\begin{itemize}
\item {\color{blue} logit} : $g(\mu_i) = g(p_i) =
\log\big(\frac{p_i}{1-p_i}\big)$ (lien canonique)
\item {\color{blue} probit} : $g(\mu_i) = g(p_i) = \Phi^{-1}(p_i)$
\item {\color{blue} log-log complémentaire} (\textit{complementary log-log}) : $g(\mu_i) = g(p_i) = \log(-\log(1-p_i))$
\end{itemize}
\end{frame}

\begin{frame}{}
\begin{center}
\includegraphics[width = 4.5cm]{figures/binarylinks.pdf}  
\end{center}

Lequel choisir ? Difficile à tester à partir des données, car une grande quantité de
données est nécessaire pour de petites valeurs de~$p$.

{\color{red} $\bigstar$ Familles générales de fonctions de lien, voir
\citeN{Aran:1981}}\\
 {\color{red}$\bigstar$ Lien non paramétrique pour les MLG, voir \citeN{Weis:Wels:1994}}\\
{\color{red}$\bigstar$ Test d'adéquation du lien, voir \citeN{Preg:1980}}
\end{frame}

\begin{frame}{}
Les valeurs  $o_i = \frac{p_i}{1-p_i}$ sont appelées les cotes (\textit{odds}). 

\bigskip

Une cote de 2/3 signifie qu'il y a 2 chances (sur 5) de succès contre 3
chances (sur 5) d'échec.

\bigskip
Note : $p_i = \frac{o_i}{1+o_i}$.

\bigskip
Le lien logit a l'avantage de permettre l'interprétation sur l'échelle des cotes.

\end{frame}


\begin{frame}{Loi de Poisson : $Y_i \sim {\mathcal P}(\lambda_i)$}


Utilisée pour modéliser des données de dénombrement (\textit{count data}), le nombre d'occurrences d'un événement dans une
période de temps ou un espace défini. Exemples : le nombre de conditions médicales rapportées par une
personne, le nombre de cyclones tropicaux pendant une saison, le nombre d'oiseaux
observés un jour donné, le nombre de fautes d'orthographe sur la page d'un
journal, le nombre d'appels entrants en 4 heures.

\bigskip
Ici, le lien garantit que $\lambda>0$. Les choix courants sont :
$g(\mu_i) = g(\lambda_i) = \log(\lambda_i)$ (lien canonique) ou $g(\mu_i) =
\sqrt{\lambda}$. 

\end{frame}


\begin{frame}{}
\begin{center}
\includegraphics[width = 4.5cm]{figures/poissonlinks.pdf}  
\end{center}

\textbf{Remarque} : le nombre d'événements sur un total est correctement modélisé par une
loi binomiale. Mais, si les probabilités de succès sont petites et les totaux
sont grands, la loi de Poisson est une bonne approximation. (Pour un petit $p$,
$\mbox{logit}(p) = \log \left(\frac{p}{1-p} \right) \simeq \log(p)$.)
\end{frame}


\begin{frame}{Loi gamma : $Y_i \sim \Gamma(\mu_i,\nu)$}

Utilisée pour modéliser des comportements continus asymétriques. Par exemple, le revenu ou le coût d'un
séjour hospitalier.

\bigskip 
Ici, le lien garantit que $\mu_i>0$. Un choix courant est :
$g(\mu_i) = \log(\mu_i)$. Le lien canonique est $1/\mu_i$ mais est plus rarement
utilisé. On note que le lien canonique ne garantit pas $\mu_i>0$.

\bigskip

\textbf{Remarque} : la loi exponentielle ${\mathcal E}(\lambda_i)$ est un
cas particulier où $\nu=1$.

\end{frame}


\begin{frame}{Extensions des MLG}

Dans l'esprit des MLG, certains modèles proches peuvent être considérés.

\bigskip

{\color{blue}Loi multinomiale :}

\bigskip
Généralisation de la loi binomiale où il y a $K$ résultats
possibles (au lieu de deux), avec des probabilités associées $\pi_1, \ldots,
\pi_K$. On note $Y_i$ la variable aléatoire prenant des valeurs parmi les catégories
$1,\ldots,K$.


\bigskip
Il faut distinguer entre les données {\color{blue}nominales} (pas d'ordre naturel) et 
les données {\color{blue}ordinales}.
\end{frame}


\begin{frame}{}
\textbf{Données nominales}\\
Exemples : couleur des yeux, affiliation politique.

\bigskip
On note $Y_{ik}$ le nombre d'observations tombant dans la catégorie $k$ pour
l'individu $i$ (si une seule mesure, alors un seul de $Y_{i1},\ldots,Y_{iK}$ est
égal à un et les autres sont zéro). Les probabilités associées sont $p_{ik} =
P(Y_i=k)$. 

\bigskip
On lie les prédicteurs $x_i$ aux probabilités par :
$$\log\left(\frac{p_{ik}}{p_{i1}} \right) = x_i^T \beta_k, \ \ \ k=2,\ldots,K$$
sous la contrainte que $\sum_{k=1}^K p_{ik}=1$, c'est-à-dire $p_{i1}=1-
\sum_{k=2}^K p_{ik}$. 


\end{frame}


\begin{frame}{}

\textbf{Données ordinales}\\
Exemples : tests de dégustation, classification de radiographies, détermination du
bien-être physique et mental. 

\smallskip
Les probabilités de réponse cumulées 
$$\gamma_1 = \pi_1, \gamma_2=\pi_1+\pi_2, \ldots, \gamma_K=1,$$
c'est-à-dire $\gamma_k(x_i) = P(Y_i \leq k \mid x_i)$ sont modélisées. 

\bigskip
Par exemple (modèle à lien logit)
$$\log\left(\frac{\gamma_k(x_i)}{1-\gamma_k(x_i)} \right) = \alpha_k - x_i^T \beta$$
pour $k=1, \ldots, (K-1)$. Le signe négatif sur le prédicteur linéaire est une
convention assurant que de grandes valeurs de $x_i^T \beta$ mènent à une augmentation de
la probabilité dans les catégories de numéro plus élevé.

\bigskip
\textbf{Extension} : permettre $\beta_k$ au lieu de $\beta$.
\end{frame}


\begin{frame}{}

Ce modèle est appelé le modèle à
{\color{blue}cotes proportionnelles} (\textit{proportional-odds}) parce que le rapport de l'événement $Y \leq k$
en $x_i=x$ et $x_i = \tilde{x}$ est
$$\frac{\mbox{odds}(Y\leq k\mid x)}{\mbox{odds}(Y\leq k\mid
\tilde{x})} = \frac{\gamma_k(x)/(1-\gamma_k(x))}{\gamma_k(\tilde{x}) / 
(1-\gamma_k(\tilde{x}))} = \exp(-(x-\tilde{x})^T \beta),$$
indépendant de $\alpha_k$. En particulier, si $x_i$ est un indicateur pour 2
groupes de traitement, $T_1$ et 
$T_2$, 
$$\frac{\mbox{odds}(Y\leq k \mid T_1)}{\mbox{odds}(Y\leq k\mid
T_2)}= \exp(-\Delta),$$
où $\Delta$ mesure l'effet du traitement.

\smallskip
Alternativement, on peut utiliser un lien log-log complémentaire
$$ \log\left(-\log(1-\gamma_k(x_i))\right)= \alpha_k - x_i^T \beta,$$
ce qui produit le modèle à {\color{blue}risques proportionnels} (\textit{proportional hazard}).

\bigskip
\textbf{Note} : dans tous ces modèles, on a besoin de $\alpha_1 \leq \alpha_2 \leq
\ldots \leq \alpha_{K-1}$  pour assurer que les probabilités sont non négatives.
\end{frame}

\subsection{Estimation des paramètres}
\label{GLM.estim.section}

\begin{frame}[allowframebreaks]{Estimation des paramètres : maximum de vraisemblance (MV)}
La fonction de vraisemblance pour la {\color{blue}famille exponentielle
canonique} est
\begin{eqnarray*}
  \label{lik.exp}
  \lefteqn{L(\theta_1,\ldots, \theta_n; y_1, \ldots, y_n)= \prod_{i=1}^n \exp \Big( A \ \frac{y_i \theta_i - b(\theta_i)}{\phi} +
c(y_i,\frac{\phi}{A})\Big) }\\
& = & \exp \Big( \frac{A}{\phi} \sum_{i=1}^n y_i \theta_i - \frac{A}{\phi}
\sum_{i=1}^n b(\theta_i) + \sum_{i=1}^n c(y_i,\frac{\phi}{A}) \Big)
\end{eqnarray*}
et la fonction de log-vraisemblance est donnée par
$$l(\theta_1,\ldots, \theta_n; y_1, \ldots, y_n)
 =  \frac{A}{\phi} \sum_{i=1}^n y_i \theta_i - \frac{A}{\phi}
\sum_{i=1}^n b(\theta_i) + \sum_{i=1}^n c(y_i,\frac{\phi}{A}).$$

La relation  $g(\mu_i) = x_i^T \beta$ est utilisée dans la
fonction de vraisemblance ci-dessus. Ensuite, la fonction est maximisée par rapport à
$\beta$ (en résolvant les équations de score) pour obtenir $\hat{\beta}$.

\bigskip

Aucune expression analytique en forme fermée pour $\hat{\beta}$ n'est disponible en général.
\end{frame}

\begin{frame}{}

\textbf{Exemple} : si $Y_i$ est une réponse binaire, c'est-à-dire $Y_i \sim
Bernoulli(p_i)$, la vraisemblance est

\begin{equation*}
  \label{lik.bernoulli}
  L(p_1,\ldots, p_n; y_1, \ldots, y_n) = \prod_{i=1}^n p_i^{y_i} (1-p_i)^{1-y_i},
\end{equation*}
et la log-vraisemblance est

\begin{equation}
  \label{loglik.bernoulli}
l(p_1,\ldots, p_n; y_1, \ldots, y_n) = \sum_{i=1}^n y_i \log (p_i) + \sum_{i=1}^n (1-y_i) \log(1-p_i). 
\end{equation}

\bigskip
Le lien logit implique que $\log \big(\frac{p_i}{1-p_i}\big) = x_i^T
\beta$, et donc que
$$p_i = \frac{\exp(x_i^T \beta)}{(1+\exp(x_i^T \beta))},$$
\end{frame}

\begin{frame}{}
ce qui donne la log-vraisemblance
\begin{eqnarray*}
 \lefteqn{l(\beta;y)=l(\beta_0, \beta_1,\ldots,\beta_p; y_1, \ldots, y_n)=}\\
& = &  \sum_{i=1}^n y_i \log \Big(\frac{\exp(x_i^T \beta)}{1+\exp(x_i^T
  \beta)} \Big) +  \sum_{i=1}^n (1-y_i) \log\Big(\frac{1}{1+\exp(x_i^T \beta)} \Big) \\
& = &  \sum_{i=1}^n y_i x_i^T \beta -  \sum_{i=1}^n y_i \log( 1 +
\exp(x_i^T \beta)) \\
& & - \sum_{i=1}^n (1-y_i) \log(1+\exp(x_i^T \beta))\\
& = & \sum_{i=1}^n y_i x_i^T \beta - \sum_{i=1}^n \log(1+\exp(x_i^T \beta))
\end{eqnarray*}
\end{frame}

\begin{frame}{}

Les équations de score (dérivation par rapport à $\beta$) sont :
$$\sum_{i=1}^n \Big(y_i  - \frac{\exp(x_i^T \beta)}{1+\exp(x_i^T
  \beta)} \Big)   x_i = 0,$$
qui doivent être résolues numériquement.

\end{frame}

\subsection{Méthodes numériques}
\begin{frame}{Procédure de Newton-Raphson dans le cas univarié}

Pour $x$ univarié, on souhaite résoudre $h(x)=0.$

\bigskip

Pour de petites distances $x^{(m)} - x^{(m-1)}$, la pente est
\begin{equation}
  \label{NR.eq}
 h^{\prime}(x^{(m-1)}) = \frac{h(x^{(m)})-h(x^{(m-1)})}{x^{(m)} - x^{(m-1)}}.
\end{equation}

\bigskip
Si $x^{(m)}$ est la solution recherchée de sorte que $h(x^{(m)})=0$, alors
\eqref{NR.eq} peut être réarrangé pour donner
\begin{equation}
  \label{NRformula}
  x^{(m)} = x^{(m-1)} - \frac{h(x^{(m-1)})}{h^{\prime}(x^{(m-1)})}.
\end{equation}
C'est la formule de Newton-Raphson pour résoudre $h(x)=0$ de manière itérative.
\label{Ufct}
\end{frame}

\begin{frame}{}
Si on applique \eqref{NRformula} pour maximiser une log-vraisemblance $l(\beta;y)$, on a
$$\beta^{(m)} = \beta^{(m-1)} - \frac{U^{(m-1)}}{U^{\prime (m-1)}},$$
où $U=  \frac{d l(\beta;y)}{d \beta}$.

\bigskip
Pour l'estimation par maximum de vraisemblance, il est courant d'approximer
$U^{\prime}$ par $E(U^{\prime}) = - {\mathcal J}$ (l'information de Fisher), pour obtenir 
\begin{equation}
  \label{scoringformula}
  \beta^{(m)} = \beta^{(m-1)} + \frac{U^{(m-1)}}{{\mathcal J}^{(m-1)}}.
\end{equation}
C'est ce qu'on appelle la {\color{blue}méthode de scoring} (\textit{method of scoring}).
\end{frame}

\begin{frame}{}
Les équivalents multivariés de  \eqref{NRformula} et
\eqref{scoringformula} sont
$$\beta^{(m)} = \beta^{(m-1)} - \big[U^{\prime (m-1)}\big]^{-1} \ U^{(m-1)},$$
et respectivement
\begin{equation}
  \label{multiscoring}
\beta^{(m)} = \beta^{(m-1)} + \big[{\mathcal J}^{(m-1)}\big]^{-1} \ U^{(m-1)},
\end{equation}
où $\beta$ et $U$ sont maintenant des vecteurs de dimension $(p+1)$ et
${\mathcal J}$ est une matrice $(p+1) \times (p+1)$.
\end{frame}

\subsection{Algorithme MCPI pour les MLG}

\begin{frame}{Algorithme de moindres carrés pondérés itératifs (MCPI)}

Étapes principales de l'algorithme de {\color{blue} moindres carrés pondérés itératifs (MCPI)}
(\textit{Iterative Weighted Least Squares}, IWLS) pour maximiser la log-vraisemblance des MLG (voir \citeN{McCu:Neld:1989},
p.\ 40-42 pour des détails supplémentaires) 

\bigskip
À résoudre :
$$\frac{\partial l(\beta;y)}{\partial \beta} =  \left(
\begin{array}{c}
\partial l(\beta;y)/\partial \beta_0\\
\vdots\\
\partial l(\beta;y)/\partial \beta_p
\end{array} \right) = 0$$

On a que
\begin{equation}
\label{scorej}
\frac{\partial l}{\partial \beta_j} = U_j = \sum_{i=1}^n
  \frac{\partial l_i}{\partial
  \beta_j} = \sum_{i=1}^n \frac{\partial l_i}{\partial \theta_i} \cdot
  \frac{\partial \theta_i}{\partial \mu_i} \cdot \frac{\partial
  \mu_i}{\partial \beta_j}.
\end{equation}
\end{frame}

\begin{frame}{}
Premièrement
$$\frac{\partial l_i}{\partial \theta_i} = A \frac{y_i -
b^{\prime}(\theta_i)}{\phi} \stackrel{\eqref{expect.exponential}}{=}
 A \frac{y_i-\mu_i}{\phi}.$$

\bigskip

Deuxièmement $\frac{\partial \theta_i}{\partial \mu_i} = 1/(\frac{\partial
  \mu_i}{\partial \theta_i})$
et
$$\frac{\partial \mu_i}{\partial \theta_i}
\stackrel{\eqref{expect.exponential}}{=} b^{\prime \prime}(\theta_i)
\stackrel{\eqref{var.exponential}}{=} A \frac{Var(Y_i)}{\phi}$$

\bigskip
Finalement
$$\frac{\partial \mu_i}{\partial \beta_j}=
\frac{\partial \mu_i}{\partial \eta_i} \cdot \frac{\partial
  \eta_i}{\partial\beta_j} = \frac{\partial \mu_i}{\partial \eta_i}
x_{ij}.$$

\bigskip

En mettant tout ensemble, \eqref{scorej} peut s'écrire
$$\frac{\partial l}{\partial \beta_j} = U_j = \sum_{i=1}^n
\frac{y_i-\mu_i}{Var(Y_i)} x_{ij} \Big(\frac{\partial \mu_i}{\partial
  \eta_i}\Big)$$
\label{EE.GLM}
\end{frame}

\begin{frame}{}
On utilise la méthode de scoring (voir~\eqref{multiscoring} à la page~\pageref{multiscoring})  ce qui implique
le calcul de ${\mathcal J}$, dont
l'élément $(j,k)$ est ${\mathcal J}_{jk} = E(U_j U_k)$ et est donné par
\begin{eqnarray*}
  \lefteqn{{\mathcal J}_{jk} = E \left\{ \sum_{i=1}^n
\frac{y_i-\mu_i}{Var(Y_i)} x_{ij} \Big(\frac{\partial \mu_i}{\partial
  \eta_i}\Big) \cdot \sum_{l=1}^n
\frac{y_l-\mu_l}{Var(Y_l)} x_{lk} \Big(\frac{\partial \mu_l}{\partial
  \eta_l}\Big) \right\}} \hspace*{3cm}\\
& = & \ldots (\mbox{quelques calculs supplémentaires})\\
& = &  \sum_{i=1}^n \frac{x_{ij} x_{ik}}{Var(Y_i)}
\Big(\frac{\partial \mu_i}{\partial
  \eta_i}\Big)^2
\end{eqnarray*}

${\mathcal J}$ peut s'écrire $X^T W X$, où $W$ est une matrice diagonale $n \times
n$ avec les éléments
$$w_{ii} = \frac{1}{Var(Y_i)} \Big(\frac{\partial \mu_i}{\partial
  \eta_i}\Big)^2$$
\end{frame}

\begin{frame}{}
La méthode de scoring est
$$\beta^{(m)} = \beta^{(m-1)} + \big[{\mathcal J}^{(m-1)}\big]^{-1} \ U^{(m-1)},  $$
et implique
$${\mathcal J}^{(m-1)} \beta^{(m)} = {\mathcal J}^{(m-1)}
\beta^{(m-1)} + U^{(m-1)}.$$

Le membre de droite de l'équation ci-dessus peut s'écrire dans notre cas
(calculs omis) comme $X^T W^{(m-1)}
z^{(m-1)}$, où $z^{(m-1)}$ est un vecteur avec les éléments
$$z_i^{(m-1)} = \sum_{k=0}^p x_{ik} \beta_k^{(m-1)} +
  (y_i-\mu_i^{(m-1)})\Big(\frac{\partial \eta_i}{\partial \mu_i}\Big)^{(m-1)},$$
où $\mu_i^{(m-1)}$ et $(\partial \mu_i/\partial \eta_i)^{(m-1)}$ sont évalués à
  $\beta^{(m-1)}$.
\end{frame}

\begin{frame}{}
La méthode de scoring revient donc à résoudre par rapport à $ \beta^{(m)}$
\begin{equation}
  \label{IWLS}
  X^T W^{(m-1)} X \beta^{(m)} = X^T W^{(m-1)} z^{(m-1)}.
\end{equation}
\label{IWLSweights}

\eqref{IWLS} est un problème de moindres carrés pondérés à résoudre
itérativement car ($z^{(m-1)}$ et $W^{(m-1)}$ dépendent de 
$\beta^{(m-1)}$). Il suffit d'un logiciel qui calcule les moindres carrés pondérés.

\bigskip
\textbf{Remarques} :
\begin{itemize}
\item $\frac{\partial \mu_i}{\partial \eta_i}$ dépend de la fonction de lien
  utilisée.
\item $Var(Y_i)$ dépend du modèle.
\item $\frac{\partial \mu_i}{\partial \eta_i}= b^{\prime \prime}(\theta) \propto Var(Y_i)$ lorsque le lien canonique
est utilisé. Ceci simplifie les équations d'estimation :
$\frac{\partial l}{\partial \beta_j} = U_j = \sum_{i=1}^n
(y_i-\mu_i) x_{ij}$
\end{itemize}
\end{frame}

\subsection{Jeu de données sur le cancer du sein}

\begin{frame}[fragile]{}

Considérons le jeu de données sur le cancer du sein de l'Exemple~2
(p.~\pageref{example2}), et ajustons le modèle simple qui explique le
résultat (\texttt{status}) en fonction
de \texttt{nuclei} seulement (plus l'ordonnée à l'origine) avec un lien logit.

\bigskip

{\scriptsize
\begin{verbatim}
Call: glm(formula = status ~ nuclei, family = binomial, 
    data = breast,  na.action = na.omit)

Coefficients:
            Estimate Std. Error z value Pr(>|z|)    
(Intercept) -3.52218    0.23205  -15.18   <2e-16 ***
nuclei       0.85935    0.07092   12.12   <2e-16 ***
---
Signif. codes:  0 '***' 0.001 '**' 0.01 '*' 0.05 '.' 0.1 ' ' 1

(Dispersion parameter for binomial family taken to be 1)

    Null deviance: 884.35  on 682  degrees of freedom
Residual deviance: 340.63  on 681  degrees of freedom
  (16 observations deleted due to missingness)
AIC: 344.63

Number of Fisher Scoring iterations: 6
\end{verbatim}
}
\end{frame}

\begin{frame}{}
Si $p_i = E(\texttt{status})$, le modèle est
$$\log\Big(\frac{p}{1-p}\Big) = \beta_0 + \beta_1 \texttt{nuclei}$$

\bigskip

\textbf{Interprétation}
\begin{itemize}
\item $\beta_1$ est relié à la pente de la courbe logistique (p.~\ref{slopelogit.fig}).
\item L'effet d'un changement unitaire de \texttt{nuclei} est d'augmenter le log des cotes de
$\beta_1$. De manière équivalente, l'effet d'un changement unitaire de  \texttt{nuclei} est de
multiplier les cotes d'une réponse positive par $\exp(\beta_1)$.
\item Lorsqu'il y a plus d'une
covariable, même interprétation à condition que les
autres covariables soient gardées fixes.
\item L'effet d'un changement unitaire de \texttt{nuclei} sur $p_i$ est plus
compliqué à établir. En présence de plusieurs variables explicatives, cet
effet dépendrait aussi des autres covariables.
\end{itemize}

\end{frame}

\begin{frame}{Courbes logistiques pour différentes valeurs de $\beta_1$}

\label{slopelogit.fig}
\begin{center}
\includegraphics[width = 7.5cm]{figures/slopelogit.pdf}  
\end{center}

\end{frame}


\begin{frame}[fragile]{}
On peut aussi ajuster le même modèle, mais avec un lien probit.


\bigskip
{\scriptsize
\begin{verbatim}
Call: glm(formula = status ~ nuclei, family = binomial(link = probit), 
    data = breast, na.action = na.omit)

Coefficients:
            Estimate Std. Error z value Pr(>|z|)    
(Intercept) -1.96960    0.10934  -18.01   <2e-16 ***
nuclei       0.45172    0.03177   14.22   <2e-16 ***
---
Signif. codes:  0 '***' 0.001 '**' 0.01 '*' 0.05 '.' 0.1 ' ' 1

(Dispersion parameter for binomial family taken to be 1)

    Null deviance: 884.35  on 682  degrees of freedom
Residual deviance: 342.94  on 681  degrees of freedom
  (16 observations deleted due to missingness)
AIC: 346.94

Number of Fisher Scoring iterations: 7
\end{verbatim}
}
\end{frame}


\begin{frame}{Logistique vs probit}

Les estimations des paramètres du modèle probit sont assez différentes de celles du
modèle logit, car elles n'ont pas la même signification :

$$\log\big(\frac{p_i}{1-p_i}\big) = x_i^T \beta^L \ \ \ \mbox{versus} \ \ \
\Phi^{-1}(p_i) = x_i^T \beta^P.$$

\bigskip

Les valeurs ajustées (prédictions pour les observations dans l'échantillon) sont
obtenues à partir de la formulation du modèle.
\end{frame}

\begin{frame}{}
Pour le modèle logistique, $\log(p_i/(1-p_i))
= x_i^T \beta^L$ implique $p_i = \exp(x_i^T \beta^L)/(1+\exp(x_i^T \beta^L))$ et
donc
$$\hat{y_i} = \hat{p_i} =  \frac{\exp(x_i^T \hat{\beta}^L)}{(1+\exp(x_i^T
\hat{\beta}^L))}.$$

\bigskip
Pour le modèle probit, $\Phi^{-1}(p_i) = x_i^T \beta^P$ implique $p_i = \Phi(x_i^T
\beta^P)$ et donc
$$ \hat{y_i} = \hat{p_i} = \Phi(x_i^T \hat{\beta}^P).$$
\end{frame}


\begin{frame}[fragile]{Valeurs ajustées avec les modèles logit et probit}

\begin{center}
\includegraphics[width = 7cm]{figures/breastsub.pdf}  
\end{center}
\end{frame}


\subsection{Vérification du modèle}

\begin{frame}{Validation du modèle}

\label{deviance.def}
Une façon d'évaluer l'adéquation d'un modèle est de le comparer à un modèle plus général
avec le nombre maximal de paramètres : un  {\color{blue}
modèle saturé}, qui a le même nombre de paramètres que
d'observations et mène à un ajustement parfait. La log-vraisemblance correspondante est 
notée $l(y;y)$.

\bigskip
La {\color{blue} déviance réduite} (\textit{scaled deviance}) est définie par :
\begin{eqnarray*}
\lefteqn{D^*(y,\mu) = 2 \big[ l(y;y) - l(\hat{\beta};y) \big]=}\\
& = & \sum_{i=1}^n 2 \big[ l_i(y_i;y_i) - l_i(\hat{\beta};y_i) \big] =
\sum_{i=1}^n d_i,
\end{eqnarray*}
et la  {\color{blue} déviance} (\textit{deviance}) $D(y,\mu) = \phi D^*(y,\mu)$.
\label{deviance.comp}
\end{frame}

\begin{frame}{}
La déviance peut être considérée comme la statistique du rapport de vraisemblance pour tester un
modèle spécifique au sein du modèle saturé, en supposant $\phi=1$. 
Cette hypothèse est
vraie pour les modèles binomiaux et de Poisson, pour  lesquels $D(y,\mu) =  D^*(y,\mu)$. Pour
d'autres distributions, par ex.\ normale ou gamma, la déviance n'est pas directement liée
à une statistique de rapport de vraisemblance. 

\bigskip
De grandes valeurs de $D$ indiquent que le modèle n'est pas bon. D'un autre
côté, de petites valeurs de $D$ (ou $D^*$) surviennent lorsque la log-vraisemblance $l(\beta;y)$
(d'un modèle à $p$ paramètres) est proche de la log-vraisemblance
$l(y;y)$ (équivalent à un modèle à $n$ paramètres).

\end{frame}


\begin{frame}{}
On a les résultats suivants :
\begin{itemize}
\item La distribution de $D^*$ est exactement
$\chi^2_{(n-(p+1))}$ si $Y_i$ suit une loi normale ($D^*$ est la somme des carrés résid.\ 
divisée par $\sigma^2$). 
\item Pour certaines autres distributions de $Y_i$, la loi $\chi^2_{(n-(p+1))}$ peut être considérée
comme une approximation de $D^*$.
\item $D$ (=$D^*$) est inutilisable lorsque $Y_i$ est Bernoulli$(p_i)$. Dans ce cas, la
distribution de $D$ est dégénérée (en fait, $D$ ne dépend pas de $Y_i$).
\item L'approximation $\chi^2_{(n-(p+1))}$ fonctionne mal pour $Y_i \sim {\mathcal
B}(m_i,p_i)$ si $m_i$ est petit (par ex.\ $<5$).
\item Sauf dans le cas Bernoulli, une procédure de
  bootstrap pourrait être
utilisée pour approximer la distribution de $D^*$.
\item Dans le cas de Bernoulli, il existe un correctif ad-hoc, voir {\color{red}
$\bigstar$ \citeN{Hosm:Leme:1980}}.
\end{itemize}
\end{frame}


\begin{frame}{}
Une statistique alternative d'adéquation est la statistique $X^2$ de Pearson :
$$X^2 = \sum_{i=1}^n \frac{(y_i - \hat{\mu}_i)^2}{v(\hat{\mu}_i)}.$$

\bigskip
La distribution de $X^2$ est $\chi^2_{n-(p+1)}$ pour des $Y_i$ normalement distribués et
cette approximation peut être utilisée pour d'autres distributions (voir la discussion sur
$D$).

\bigskip

On peut démontrer que $D$ et $X^2$ sont reliés.
\end{frame}


\begin{frame}{Mesure globale de la qualité de l'ajustement}

\begin{itemize}
\item {\color{blue}Proportion de déviance expliquée}
$$(D_{null}-D)/D_{null},$$
où $D_{null}$ est la déviance de l'ajustement avec l'ordonnée à l'origine seulement.
\item Une {\color{blue}extension de la mesure $R^2$} \cite{Nage:1991} :
\begin{eqnarray*}
\widetilde{R}^2 & = & \frac{1-(L_{null}/L)^{2/n}}{1-L_{null}^{2/n}}\\
\end{eqnarray*}
où $L_{null}$ est la vraisemblance de l'ajustement avec l'ordonnée à l'origine seulement.
\end{itemize}
\end{frame}

\begin{frame}{Analyse des résidus}

Plusieurs types de résidus peuvent être définis pour les modèles MLG. En particulier, 
\begin{itemize}
\item {\color{blue}résidus de Pearson} : $r_{iP} = \frac{y_i - \hat{\mu}_i}{\sqrt{v(\hat{\mu}_i)}},$
\item {\color{blue}résidus de déviance} (\textit{deviance residuals}) : $r_{iD} = \mbox{sign}(y_i - \hat{\mu_i}) \sqrt{d_i},$
où $d_i$ sont les composantes individuelles de la déviance $D$ (voir page~\pageref{deviance.comp}).
\item {\color{blue}résidus de déviance standardisés} : $r_{iDS} = \frac{r_{iD}}{\sqrt{1-h_{ii}}},$
où $h_{ii}$ sont les éléments diagonaux de la matrice chapeau (\textit{hat matrix}) $H=X(X^T
X)^{-1}X^T$.
\end{itemize}

\bigskip
\textbf{Note} : \\
$\sum_{i=1}^n r_{iP}^2 = X^2$ et $\sum_{i=1}^n r_{iD}^2 =
\sum_{i=1}^n d_i = D$.

\end{frame}

\begin{frame}{}
Les résidus sont habituellement représentés graphiquement :
\begin{itemize}
\item en fonction des valeurs ajustées pour vérifier le manque d'ajustement (écarts/structures).
\item en fonction des quantiles normaux (diagramme QQ). 
\item dans l'ordre des mesures pour vérifier la corrélation sérielle.
\item en fonction de chaque variable explicative continue pour vérifier l'hypothèse de linéarité.
\end{itemize}

\bigskip
\textbf{Remarques} : 
\begin{itemize}
\item Des structures apparaissent très souvent pour $Y_i \sim {\mathcal B}(m_i,p_i)$ pour de petites valeurs de $m_i$.
\item Les facteurs dans les covariables peuvent induire des structures particulières dans les résidus.
\end{itemize}

\bigskip
\textbf{Note} : 
La normalité des résidus n'est pas attendue, et donc le
diagramme QQ normal doit être interprété avec prudence.
\end{frame}

\begin{frame}{}
Illustration avec des données simulées à partir du modèle 
$$Y_i \sim Bernoulli(p_i),$$
avec $p_i = \exp(x_i^T \beta) /(1+ \exp(x_i^T \beta))$, où $x_i^T = (1, z_i)$ avec $z_i \sim N(0,1)$ et $\beta^T = (1, 0.5)$.

\begin{center}
\includegraphics[width = 5.5cm]{figures/QQnormBernoulliSim.pdf}  
\end{center}

\end{frame}

\begin{frame}{Alternative : résidus quantiles randomisés}

Soit $F = F_{\theta_i}$ la fonction de répartition de $Y_i$ et
$\Phi$ la fonction de répartition d'une variable aléatoire normale standard.

\smallskip
Les « résidus quantiles randomisés » (\textit{randomized quantile residuals}) sont définis par :
\begin{itemize}
\item si $F$ est continue : $r_{iQ} = \Phi^{-1}(F(y_i; \hat{\mu}_i,\hat{\phi}))$. 
\item si $F$ est discrète : $r_{iQ} = \Phi^{-1}(u_i)$, avec $u_i$ une variable aléatoire uniforme sur l'intervalle $(a_i,b_i]$, avec $a_i = \lim_{y \rightarrow y_i} F(y;
\hat{\mu}_i,\hat{\phi})$ et $b_i=F(y_i; \hat{\mu}_i,\hat{\phi})$, 
\end{itemize}

La distribution des « résidus quantiles randomisés » est exactement
normale standard, à la variabilité d'échantillonnage près de $\hat{\mu}_i$ et $ \hat{\phi}$ (grâce à la transformation intégrale de probabilité).

\smallskip
Note : Lorsque $Y_i$ est discret, la définition inclut une étape de randomisation : utile de regarder plusieurs réalisations.
\end{frame}


\begin{frame}[fragile]{}
Sur l'exemple logistique simulé précédent :

\begin{center}
\includegraphics[width = 5.5cm]{figures/QQrandomizedBernoulliSim.pdf}  
\end{center}

{\scriptsize
\begin{verbatim}
> require(statmod) 
> qqnorm(qresid(glm.sim))
\end{verbatim}
}
\end{frame}

\subsection{Exemple 4 : Gamma}

\begin{frame}[fragile]{Exemple 4 de la page~\pageref{CHUV.Ex}} 
\label{Routput.gamma}


\bigskip 
{\tiny
\begin{verbatim}
Call: glm(formula = CouTot ~ log(LOS) + Typadm + Typass + age + Sexe + 
    dest, family = Gamma(link = log), data = MYdata)

Coefficients:
              Estimate Std. Error t value Pr(>|t|)    
(Intercept)  7.2338121  0.1468364  49.264  < 2e-16 ***
log(LOS)     0.8222203  0.0279641  29.403  < 2e-16 ***
Typadm       0.2136148  0.0500028   4.272 4.67e-05 ***
Typass       0.0932763  0.0791100   1.179   0.2414    
age         -0.0005335  0.0012852  -0.415   0.6790    
Sexe         0.0951009  0.0499814   1.903   0.0602 .  
dest        -0.1043560  0.0692713  -1.506   0.1353    
---
Signif. codes:  0 '***' 0.001 '**' 0.01 '*' 0.05 '.' 0.1 ' ' 1

(Dispersion parameter for Gamma family taken to be 0.04961642)

    Null deviance: 47.9536  on 99  degrees of freedom
Residual deviance:  5.0718  on 93  degrees of freedom
AIC: 1817.9

Number of Fisher Scoring iterations: 5
\end{verbatim}
}
\end{frame}


\begin{frame}{}
La déviance est $D=5.07$ dans cet exemple, voir \texttt{Residual
  Deviance} dans \texttt{summary(MYglm)}.

\bigskip
On l'utilise pour tester l'hypothèse $H_0$ que notre modèle ajuste bien les données.

\bigskip
La valeur-$p$ est $P(D > 5.07)$, qui doit être calculée selon une
loi $\chi^2_{93}$. Cette probabilité est très proche de 1, ne fournissant aucune
preuve contre $H_0$.

\end{frame}

\begin{frame}[fragile]{Résidus de Pearson et de déviance}

\begin{center}
\includegraphics[width = 7cm]{figures/MYresid.pdf}  
\end{center}
\end{frame}



\begin{frame}{Résidus de Pearson en fonction des covariables continues}
\begin{center}
\includegraphics[width = 7cm]{figures/MYresid1.pdf}  
\end{center}
\end{frame}



\begin{frame}{}

\begin{center}
\includegraphics[width = 5.5cm]{figures/QQresid.pdf}  
\end{center}

\textbf{Note} :  \texttt{plot()} dans R produit
des graphiques similaires, plus le graphique des distances de Cook qui permet d'identifier
les observations influentes.
\end{frame}



\begin{frame}{}
\begin{center}
\includegraphics[width = 5cm]{figures/QQresidQrandom.pdf}  
\end{center}

Les graphiques sont identiques ! (pas de randomisation impliquée).

\bigskip
Il semble y avoir quelques valeurs aberrantes ; ce jeu de données sera réanalysé au
Chapitre~\ref{robust.section}.
\end{frame}

\subsection{Intervalles de confiance et tests z}
\begin{frame}{Intervalles de confiance}
L'estimateur de $\beta$ dans un MLG est habituellement obtenu par maximum
de vraisemblance. Cela signifie que $\hat{\beta}$ hérite des propriétés générales
du maximum de vraisemblance, en particulier la distribution
asymptotique. À partir de
$$ (\hat{\beta}- \beta) \sim {\mathcal N} (0,\tilde{\mathcal J}^{-1}), $$
où $\tilde{J}$ est la matrice d'information calculée à partir de la vraisemblance conjointe, le calcul d'une approximation
asymptotique de l'erreur standard de $\hat{\beta}$ est possible.
\end{frame}


\begin{frame}{}
Du résultat de normalité asymptotique, il découle que des intervalles de
confiance de niveau $(1-\alpha)$ pour chaque paramètre peuvent être construits comme : 
$$\big[\hat{\beta}_i - z_{1-\alpha/2} \sqrt{(\tilde{\mathcal J}^{-1})_{ii}} ;\hat{\beta}_i +
z_{1-\alpha/2} \sqrt{(\tilde{\mathcal J}^{-1})_{ii}} \big],$$
où $z_{1-\alpha/2}$ est le quantile $(1-\alpha/2)$ d'une loi ${\mathcal
  N}(0,1)$.

\bigskip
Ces intervalles sont symétriques par construction.

\bigskip
Ces intervalles de confiance peuvent mal fonctionner, en particulier si la distribution de l'estimateur du paramètre 
est nettement asymétrique. 
\end{frame}


\begin{frame}{}
En revanche, les {\color{blue}intervalles de confiance de vraisemblance profilée} (\textit{profiled likelihood confidence intervals}) ne supposent pas
la normalité de l'estimateur et fonctionnent mieux pour les petits échantillons. Ils sont
obtenus en inversant les statistiques du rapport de vraisemblance, et à ce titre ils sont
encore basés sur une approximation asymptotique (la loi $\chi^2$ de la
statistique du rapport de log-vraisemblance). 

\bigskip
Pour un modèle avec un paramètre d'intérêt $\theta$ et $\delta$ (paramètres supplémentaires),
la vraisemblance profilée est 
$$L_1(\theta) = \max_{\delta} L(\theta,\delta).$$
\end{frame}


\begin{frame}{}
Un intervalle de confiance de vraisemblance profilée de niveau $(1-\alpha)$ est l'ensemble de toutes les valeurs
$\theta_0$ telles qu'un test bilatéral de l'hypothèse nulle $H_0: \theta =
\theta_0$ ne serait pas rejeté au niveau $\alpha$. La statistique du test du rapport de vraisemblance
est utilisée.

\bigskip
Ces intervalles de confiance ne sont pas symétriques, bien que le résultat soit habituellement
proche des intervalles de confiance obtenus à partir de la théorie du maximum de vraisemblance.
\end{frame}


\begin{frame}{}
\textbf{Valeurs z} pour tester ($H_0: \beta_i=0$ contre $H_A: \beta_i \neq
0$), ce qui  correspond à un test sur la variable $x_i$ :
$$\texttt{z-value} = \frac{\hat{\beta}_i}{s.e.(\hat{\beta}_i)} =
\frac{\hat{\beta}_i}{\sqrt{({\mathcal J}^{-1})_{ii}}} \stackrel{H_0}{\sim} {\mathcal
  N}(0,1).$$
Puis on calcule la valeur-$p$
\begin{eqnarray*}
\lefteqn{P(|\texttt{z-value}| > |z_{obs}|) = \hspace*{3cm}}\\
& = & P(\texttt{z-value} < - |z_{obs}|) + P(\texttt{z-value} > |z_{obs}|) =\\
& = & 2 (1-P(\texttt{z-value} > |z_{obs}|)) =2 ( 1 - \Phi^{-1}(|z_{obs}|)).
\end{eqnarray*}

\bigskip
\textbf{Note :} Si la distribution dépend aussi d'un paramètre de
nuisance ou supplémentaire (normale, gamma), la loi $t_{n-(p+1)}$ est utilisée à la place de la
loi normale comme distribution de référence de la statistique \texttt{z} sous $H_0$. Les \texttt{z-values} sont appelées
\texttt{t-values} dans ce cas, voir par exemple la sortie de l'exemple gamma
à la page~\pageref{Routput.gamma}.
\end{frame}

\subsection{Tests d'hypothèses}
\begin{frame}{Tests d'hypothèses}
Pour tester une hypothèse paramétrique sur $\beta$, par ex.\ $H_0:
\beta=\tilde{\beta}$, les trois tests classiques sont disponibles :

\begin{description}
\item[Statistique de Wald] $(\hat{\beta} -\tilde{\beta})^T {\mathcal J}(\tilde{\beta})
  (\hat{\beta} -\tilde{\beta})$ qui suit une loi $\chi^2_{p+1}$
  sous $H_0$.
\item[Statistique du score] $U(\tilde{\beta})^T {\mathcal
J(\tilde{\beta})}^{-1}U(\tilde{\beta})$ qui suit une loi 
  $\chi^2_{p+1}$ sous $H_0$, ($U$ est la fonction de score définie à la
  page~\pageref{Ufct}).
\item[Statistique du rapport de vraisemblance] $2
  [l(\hat{\beta};y)-l(\tilde{\beta};y)]$ qui est en fait une différence
  de déviances, voir ci-dessous.
\end{description}
\end{frame}

\subsection{Sélection de variables}

\begin{frame}{Sélection de variables}

Comparer deux modèles : un ${\mathcal M}_{p+1}$ avec $(p+1)$
paramètres et un modèle {\color{blue} emboîté} ${\mathcal M}_{p+1-q}$
avec $(p+1-q)$ paramètres (c'est-à-dire un modèle contenant un sous-ensemble des
variables de ${\mathcal M}_{p+1}$).

\bigskip
La pratique courante en MLG est d'utiliser la différence de déviances pour la
sélection de variables (modèles).

\bigskip
Pour tester l'hypothèse nulle
$$H_0: \beta = (\cdot, \ldots, \cdot, 0_{p+2-q},\ldots,0_{p+1})$$
(le modèle avec $(p+1-q)$ paramètres suffit), on
définit la différence de déviances réduites
\begin{eqnarray*}
  \lefteqn{\Delta D^* = D^*(y,\hat{\mu}^{p+1-q}) - D^*(y,\hat{\mu}^{p+1}) = }\\ 
& = & 2 \Big[ l(\hat{\beta}^{p+1};y) - l(\hat{\beta}^{p+1-q};y)\Big] = \frac{D(y,\hat{\mu}^{p+1-q}) - D(y,\hat{\mu}^{p+1})}{\phi}
\end{eqnarray*}

\end{frame}



\begin{frame}{}
Si $\phi$ est connu et sous $H_0$, la distribution asymptotique de $\Delta D^*$ est
$\chi^2_{q}$ (c'est la statistique du rapport de vraisemblance). Cette approximation 
est bien meilleure que l'approximation de la déviance elle-même. 

\bigskip
Lorsque $\phi$ n'est pas connu (par ex.\ normale, gamma, modèles de type « quasi ») l'approximation
habituelle sous 
$H_0$ utilise une statistique $F$ :
$$\frac{(D(y,\hat{\mu}^{p+1-q}) - D(y,\hat{\mu}^{p+1}))/q}{\hat{\phi}} \sim
F_{q,n-(p+1)},$$
où $\hat{\phi}= \frac{D(y,\hat{\mu}^{p+1})}{n-(p+1)}$. 

\bigskip
Notons que pour le cas normal-identité, c'est un résultat exact, mais pour la
loi gamma, la précision de cette approximation n'est pas bien connue.
\end{frame}

\begin{frame}{}
La différence de déviances est souvent utilisée dans une approche séquentielle, soit
ascendante (\textit{forward}, à partir du modèle nul/constant) soit
descendante (\textit{backward}, à partir du modèle complet).

\bigskip
{\color{blue}Inconvénient} de cette approche : le modèle final dépend du
« chemin » (a priori) arbitraire de la procédure.
\end{frame}


\begin{frame}{}
Une approche complète serait d'utiliser un critère général pour comparer tous
les modèles possibles. Un tel critère est le {\color{blue}critère d'information
d'Akaike (AIC)}, qui pour un modèle général ${\mathcal M}_{d}$ avec $d$
paramètres est défini par :
\begin{eqnarray*}
\lefteqn{-2 l(\hat{\beta}^d;y) + 2d  \varpropto   D^*(y,\hat{\beta}^d) + 2d} \\
& = & \frac{D(y,\hat{\beta}^d)}{\phi} + 2 d \varpropto  D(y,\hat{\beta}^d) + 2 \phi d.
\end{eqnarray*}
Il faut estimer $\phi$.

\bigskip
Ce critère tient compte de la complexité du modèle
(mesurée par le nombre de paramètres).
Les modèles avec un AIC faible doivent être préférés.

\bigskip
Non réalisable de comparer tous les modèles possibles pour un grand $p$ ($2^p$
modèles). Utiliser l'AIC pas-à-pas (\textit{stepwise}).
\end{frame}


\subsection{Jeu de données sur le cancer du sein}
\begin{frame}[fragile]{Jeu de données sur le cancer du sein}
{\tiny
\begin{verbatim}
Coefficients:
             Estimate Std. Error z value Pr(>|z|)    
(Intercept) -10.10394    1.17488  -8.600  < 2e-16 ***
clump         0.53501    0.14202   3.767 0.000165 ***
cellsize     -0.00628    0.20908  -0.030 0.976039    
cellshape     0.32271    0.23060   1.399 0.161688    
adhesion      0.33064    0.12345   2.678 0.007400 ** 
singlesize    0.09663    0.15659   0.617 0.537159    
nuclei        0.38303    0.09384   4.082 4.47e-05 ***
chromatin     0.44719    0.17138   2.609 0.009073 ** 
nucleoli      0.21303    0.11287   1.887 0.059115 .  
mitoses       0.53484    0.32877   1.627 0.103788    
---
Signif. codes:  0 '***' 0.001 '**' 0.01 '*' 0.05 '.' 0.1 ' ' 1

(Dispersion parameter for binomial family taken to be 1)

    Null deviance: 884.35  on 682  degrees of freedom
Residual deviance: 102.89  on 673  degrees of freedom
  (16 observations deleted due to missingness)
AIC: 122.89

Number of Fisher Scoring iterations: 8
\end{verbatim}
}
\end{frame}


\begin{frame}{Résidus randomisés}
\begin{center}
\includegraphics[width = 7cm]{figures/breastQQrandom.pdf}  
\end{center}
\end{frame}


\begin{frame}[fragile]{}
Intervalles de confiance profilés à $95\%$ :

{\tiny
\begin{verbatim}
> round(confint(breast.fit,level=0.95),4)
Waiting for profiling to be done...
               2.5 %  97.5 %
(Intercept) -12.7581 -8.0851
clump         0.2742  0.8377
cellsize     -0.3948  0.4379
cellshape    -0.1482  0.7684
adhesion      0.0930  0.5869
singlesize   -0.2169  0.4049
nuclei        0.2069  0.5787
chromatin     0.1232  0.7998
nucleoli     -0.0018  0.4456
mitoses      -0.0067  1.1065
\end{verbatim}
}

et  intervalles de confiance à $95\%$ basés sur la normalité asymptotique :
{\tiny
\begin{verbatim}
> round(confint.default(breast.fit,level=0.95),4)
               2.5 %  97.5 %
(Intercept) -12.4067 -7.8012
clump         0.2567  0.8134
cellsize     -0.4161  0.4035
cellshape    -0.1293  0.7747
adhesion      0.0887  0.5726
singlesize   -0.2103  0.4036
nuclei        0.1991  0.5670
chromatin     0.1113  0.7831
nucleoli     -0.0082  0.4343
mitoses      -0.1095  1.1792
\end{verbatim}
}
\end{frame}


\begin{frame}[fragile]{Analyse de la déviance (séquentielle) :}

{\tiny
\begin{verbatim}
> anova(breast.fit,test="Chisq")

Terms added sequentially (first to last)

           Df Deviance Resid. Df Resid. Dev  Pr(>Chi)    
NULL                         682     884.35              
clump       1   425.87       681     458.48 < 2.2e-16 ***
cellsize    1   261.91       680     196.58 < 2.2e-16 ***
cellshape   1    20.08       679     176.50 7.427e-06 ***
adhesion    1    21.39       678     155.11 3.750e-06 ***
singlesize  1     6.45       677     148.66   0.01111 *  
nuclei      1    28.97       676     119.69 7.348e-08 ***
chromatin   1     9.21       675     110.48   0.00241 ** 
nucleoli    1     3.87       674     106.61   0.04906 *  
mitoses     1     3.72       673     102.89   0.05378 .  
---
Signif. codes:  0 '***' 0.001 '**' 0.01 '*' 0.05 '.' 0.1 ' ' 1
\end{verbatim}
}

Le modèle avec ordonnée à l'origine seulement est d'abord comparé au modèle
avec \texttt{clump}. La déviance pour tester $H_0:
\beta_{\texttt{clump}}=0$ passe de 884.35 à 458.48 ($\Delta D = 425.87$) (valeur-$p$, selon $\chi^2_1$, proche de 0). On rejette $H_0$.
Ensuite on teste $H_0: \beta_{\texttt{cellsize}}=0$ : $\Delta D = 261.91$ et valeur-$p$
$\simeq 0$. $H_0$ est rejeté.
\end{frame}


\begin{frame}[fragile]{}
L'analyse dépend de l'ordre dans lequel les variables sont entrées dans le modèle :

{\scriptsize
\begin{verbatim}
> breast.fit2 <- glm(status~clump+cellshape+adhesion+singlesize+
nuclei+chromatin+nucleoli+mitoses+cellsize,data=breast,
family=binomial,na.action=na.omit)

> anova(breast.fit2,test="Chisq")

Terms added sequentially (first to last)

           Df Deviance Resid. Df Resid. Dev  Pr(>Chi)    
NULL                         682     884.35              
clump       1   425.87       681     458.48 < 2.2e-16 ***
cellshape   1   258.09       680     200.39 < 2.2e-16 ***
adhesion    1    40.17       679     160.23 2.330e-10 ***
singlesize  1     9.44       678     150.79  0.002127 ** 
nuclei      1    29.38       677     121.41 5.954e-08 ***
chromatin   1    10.58       676     110.83  0.001145 ** 
nucleoli    1     4.18       675     106.66  0.041022 *  
mitoses     1     3.77       674     102.89  0.052188 .  
cellsize    1     0.00       673     102.89  0.976067    
\end{verbatim}
}
\end{frame}

\begin{frame}[fragile]{}
AIC sur tous les modèles. Ici $2^9 = 512$ modèles possibles. Lorsqu'il y a trop de modèles, peut être
utilisé de manière pas-à-pas.

{\tiny
\begin{verbatim}
> step(breast.fit)
Start:  AIC=122.89
status ~ clump + cellsize + cellshape + adhesion + singlesize + 
    nuclei + chromatin + nucleoli + mitoses

             Df Deviance    AIC
- cellsize    1   102.89 120.89
- singlesize  1   103.27 121.27
- cellshape   1   104.74 122.74
<none>            102.89 122.89
- mitoses     1   106.61 124.61
- nucleoli    1   106.66 124.66
- adhesion    1   110.31 128.31
- chromatin   1   110.33 128.33
- clump       1   120.72 138.72
- nuclei      1   122.07 140.07

\end{verbatim}
}
\end{frame}

\begin{frame}[fragile]{}
{\tiny
\begin{verbatim}
Step:  AIC=120.89
status ~ clump + cellshape + adhesion + singlesize + nuclei +  chromatin + nucleoli + mitoses

             Df Deviance    AIC
- singlesize  1   103.27 119.27
<none>            102.89 120.89
- mitoses     1   106.66 122.66
- cellshape   1   106.66 122.66
- nucleoli    1   106.76 122.76
- adhesion    1   110.64 126.64
- chromatin   1   110.70 126.70
- clump       1   121.10 137.10
- nuclei      1   122.07 138.07

Step:  AIC=119.27
status ~ clump + cellshape + adhesion + nuclei + chromatin +    nucleoli + mitoses

            Df Deviance    AIC
<none>           103.27 119.27
- mitoses    1   107.14 121.14
- nucleoli   1   107.72 121.72
- cellshape  1   107.90 121.90
- chromatin  1   111.69 125.69
- adhesion   1   112.17 126.17
- clump      1   121.55 135.55
- nuclei     1   123.15 137.15

\end{verbatim}
}
\end{frame}

\begin{frame}[fragile]{}

{\tiny
\begin{verbatim}
Call:  glm(formula = status ~ clump + cellshape + adhesion + nuclei + 
    chromatin + nucleoli + mitoses, family = binomial, data = breast, 
    na.action = na.omit)

Coefficients:
(Intercept)    clump   cellshape   adhesion       nuclei    chromatin   nucleoli    mitoses  
    -9.9828   0.5340      0.3453    0.3425       0.3883       0.4619     0.2261     0.5312  

Degrees of Freedom: 682 Total (i.e. Null);  675 Residual
  (16 observations deleted due to missingness)
Null Deviance:	    884.4 
Residual Deviance: 103.3 	AIC: 119.3
\end{verbatim}
}

\bigskip
À chaque étape, tous les sous-modèles avec une variable de moins sont considérés, et
celui avec le plus petit AIC est choisi. La procédure est arrêtée lorsqu'aucune
amélioration de l'AIC n'est observée.

\end{frame}

\begin{frame}[fragile]{}
Une recherche complète donnerait :

{\tiny
\begin{verbatim}
> require(MuMIn)
> options(na.action = "na.fail")
> breast.woNA <- na.omit(breast)
> breast.fit.woNA <- glm(status~clump+cellsize+cellshape+adhesion+singlesize+nuclei+chromatin+nucleoli+mitoses,data=breast.woNA,family=binomial)
> breast.all <- dredge(breast.fit.woNA,rank = "AIC")
> head(breast.all)

Model selection table 
    (Intrc)  adhsn  cllsh    cllsz  chrmt  clump  mitss  nucle
252  -9.983 0.3425 0.3453          0.4619 0.5340 0.5312 0.3883
508 -10.100 0.3299 0.3182          0.4465 0.5346 0.5341 0.3831
220  -9.767 0.3375 0.3495          0.4713 0.6225        0.3786
256  -9.975 0.3415 0.3396 0.007847 0.4610 0.5334 0.5302 0.3883
254 -10.010 0.3450        0.233300 0.4794 0.5759 0.5310 0.4203
124  -9.990 0.3502 0.4655          0.5203 0.5478 0.5690 0.3871  
          
     nucll   sngls    df  logLik   AIC delta weight
252 0.2261             8 -51.633 119.3  0.00  0.351
508 0.2125 0.09612     9 -51.445 120.9  1.62  0.156
220 0.2432             7 -53.572 121.1  1.88  0.137
256 0.2252             9 -51.633 121.3  2.00  0.129
254 0.2500             8 -52.677 121.4  2.09  0.124
124                    7 -53.861 121.7  2.46  0.103
Models ranked by AIC(x) 
\end{verbatim}
}

\end{frame}

\subsection{Ce qui peut mal tourner... Surdispersion}

\begin{frame}[fragile]{Navires de transport}

Dommages causés par les vagues à la
section avant de certains navires de transport de marchandises (tiré de \citeN{McCu:Neld:1989},
pages 204-208).

\bigskip
La variable réponse est le nombre d'incidents de dommages. Les variables explicatives sont le
type de navire (de $A$ à $E$), l'année de construction (classes 1960-64, 1965-69,
1970-74, 1975-79), la période d'opération (1960-74, 1975-79) et le nombre total cumulé
de mois de service.

{\scriptsize
\begin{verbatim}
> head(ship)
  Ship    Year  Period Aggregate Incidents
1    A 1960-64 1960-74       127         0
2    A 1960-64 1975-79        63         0
3    A 1965-69 1960-74      1095         3
4    A 1965-69 1975-79      1095         4
5    A 1970-74 1960-74      1512         6
6    A 1970-74 1975-79      3353        18
\end{verbatim}}
\end{frame}

\begin{frame}[fragile]{}
{\scriptsize
\begin{verbatim}
> ftable(xtabs(Incidents~Ship+Year+Period,data=ship))
             Period 1960-74 1975-79
Ship Year                          
A    1960-64              0       0
     1965-69              3       4
     1970-74              6      18
     1975-79              0      11
B    1960-64             39      29
     1965-69             58      53
     1970-74             12      44
     1975-79              0      18
C    1960-64              1       1
     1965-69              0       1
     1970-74              6       2
     1975-79              0       1
D    1960-64              0       0
     1965-69              0       0
     1970-74              2      11
     1975-79              0       4
E    1960-64              0       0
     1965-69              7       7
     1970-74              5      12
     1975-79              0       1
\end{verbatim}
}
\end{frame}

\begin{frame}{}
Le modèle de Poisson à lien log postulé est :
\begin{eqnarray*}
\lefteqn{\log(E(\texttt{Incidents})) = \hspace*{4cm}}\\
& = & \beta_0 + \log(\texttt{Aggregate}) + \beta_1 \texttt{Ship} + \beta_2
\texttt{Year} \\
 & & + \beta_3 \texttt{Period}.
\end{eqnarray*}

\textbf{Remarque} : $\beta_1$, $\beta_2$ et $\beta_3$ sont des vecteurs, dont la longueur est
égale au nombre de niveaux du facteur correspondant moins un.
\end{frame}

\begin{frame}{Modèle de taux}
Le terme
$\log(\texttt{Aggregate})$ a un coefficient $\beta_{s}=1$, connu a priori (un
{\color{blue}décalage} (\textit{offset}) dans la terminologie des MLG). Cette variable est incluse dans le modèle
pour contrôler le nombre d'occasions pour que l'événement se produise (effort).

\bigskip
On note l'équivalence entre :
\begin{eqnarray*}
\log(E(\texttt{Incidents})) = \beta_0 + \log(\texttt{Aggregate}) + \cdots \\
\end{eqnarray*}
et
\begin{eqnarray*}
\log(E(\texttt{Incidents})/\texttt{Aggregate}) =  \beta_0 + \cdots 
\end{eqnarray*}
\end{frame}

\begin{frame}{}
Dans des cas comme celui ci-dessus, où la variable d'effort est un compte, la 
modélisation peut être vue comme une approximation de Poisson à un modèle binomial.  Dans d'autres
situations (la variable d'effort est le temps, par exemple) cela n'est pas possible.

\bigskip
\bigskip
 La pertinence de fixer $\beta_{s}=1$ peut être vérifiée en ajustant un modèle où
 $\beta_{s}$ est également estimé à partir des données. 
\end{frame}

\begin{frame}[fragile]{}

{\scriptsize
\begin{verbatim}
> ship.glm <- glm(Incidents~Ship+Year+Period+offset(log(Aggregate)),
family=poisson,data=ship)
> summary(ship.glm)

Coefficients:
              Estimate Std. Error z value Pr(>|z|)    
(Intercept)   -6.40590    0.21744 -29.460  < 2e-16 ***
ShipB         -0.54334    0.17759  -3.060  0.00222 ** 
ShipC         -0.68740    0.32904  -2.089  0.03670 *  
ShipD         -0.07596    0.29058  -0.261  0.79377    
ShipE          0.32558    0.23588   1.380  0.16750    
Year1965-69    0.69714    0.14964   4.659 3.18e-06 ***
Year1970-74    0.81843    0.16977   4.821 1.43e-06 ***
Year1975-79    0.45343    0.23317   1.945  0.05182 .  
Period1975-79  0.38447    0.11827   3.251  0.00115 ** 
---
Signif. codes:  0 '***' 0.001 '**' 0.01 '*' 0.05 '.' 0.1 ' ' 1

(Dispersion parameter for poisson family taken to be 1)

    Null deviance: 146.328  on 33  degrees of freedom
Residual deviance:  38.695  on 25  degrees of freedom
AIC: 154.56
\end{verbatim}
}

\end{frame}

\begin{frame}{}
Test d'adéquation avec la statistique de déviance :
$$P(D>38.695) =  0.0395,$$
où on a utilisé une loi $\chi^2_{25}$ pour $D$.

On rejette donc (au niveau $5\%$) l'hypothèse nulle que le
modèle ajuste bien les données.

\begin{center}
\includegraphics[width = 5cm]{figures/shipdiag.pdf}  
\end{center}
\end{frame}

\begin{frame}{Résidus randomisés}
\begin{center}
\includegraphics[width = 5.5cm]{figures/shipQQrandom.pdf}  
\end{center}

{\color{blue}Raisons du manque d'ajustement ?}

\end{frame}


\begin{frame}{}
Si $Var(Y_i) > E(Y_i)$, le modèle de Poisson ne tient plus et ce
phénomène est appelé {\color{blue}surdispersion} (\textit{overdispersion}). (Si $Var(Y_i) < E(Y_i)$, on
a une {\color{blue}sous-dispersion} (\textit{underdispersion}) mais celle-ci est moins courante).

\bigskip
L'une des principales raisons de la surdispersion est le regroupement dans la population
(le paramètre varie d'un groupe à l'autre, en fonction de la taille du groupe par
exemple). D'un point de vue statistique, le paramètre est considéré comme aléatoire plutôt
que fixe.

\bigskip
Une grande valeur de la déviance peut indiquer une surdispersion.

\bigskip
Une première vérification sur les données est de supposer que $Var(Y_i) = \sigma^2
v_i = \sigma^2 \mu_i$ et de regarder la valeur ajustée de $\sigma^2$. C'est un
ajustement \textit{ad-hoc} et ne correspond pas nécessairement à une
vraisemblance.
\end{frame}

\begin{frame}{}
Le modèle ne peut plus être ajusté par maximum de vraisemblance, mais par un
ensemble d'équations d'estimation différentes :
$$\sum_{i=1}^n \Big(\frac{y_i-\mu_i}{\sigma^2 v(\mu_i)}\Big) \ \frac{\partial
  \mu_i}{\partial \beta}=0,$$
qui correspondent à des fonctions de quasi-vraisemblance (\textit{quasi-likelihood})
$$\sum_{i=1}^n \int_{y_i}^{\mu_i} \frac{y_i-t}{\sigma^2 v(t)} dt.$$
\label{QL.page}
\end{frame}

\begin{frame}{}
\textbf{Remarques} :
\begin{itemize}
\item Les approches par maximum de vraisemblance et par quasi-vraisemblance donnent les mêmes
estimateurs pour tous les modèles de la famille exponentielle à un paramètre (binomial,
Poisson, par exemple).
\item Notons que $\sigma^2$ n'affecte pas les équations de score, mais
  a un impact multiplicatif sur l'erreur standard des coefficients (et sur la 
  distribution de la différence de déviances).
\end{itemize}

\bigskip
{voir \color{red} $\bigstar$ \citeN{Gani:Scha:1992}, \citeN{Lamb:Roed:1995} et \citeN{Dean:Lawl:1989}}
\end{frame}

\begin{frame}[fragile]{}
\label{shipmodel.page}

{\scriptsize
\begin{verbatim}
> ship.glm.quasi <- glm(Incidents~Ship+Year+Period+offset(log(Aggregate)),
family=quasipoisson,data=ship)
> summary(ship.glm.quasi)

Coefficients:
              Estimate Std. Error t value Pr(>|t|)    
(Intercept)   -6.40590    0.28276 -22.655  < 2e-16 ***
ShipB         -0.54334    0.23094  -2.353  0.02681 *  
ShipC         -0.68740    0.42789  -1.607  0.12072    
ShipD         -0.07596    0.37787  -0.201  0.84230    
ShipE          0.32558    0.30674   1.061  0.29864    
Year1965-69    0.69714    0.19459   3.583  0.00143 ** 
Year1970-74    0.81843    0.22077   3.707  0.00105 ** 
Year1975-79    0.45343    0.30321   1.495  0.14733    
Period1975-79  0.38447    0.15380   2.500  0.01935 *  
---
Signif. codes:  0 '***' 0.001 '**' 0.01 '*' 0.05 '.' 0.1 ' ' 1

(Dispersion parameter for quasipoisson family taken to be 1.691028)

    Null deviance: 146.328  on 33  degrees of freedom
Residual deviance:  38.695  on 25  degrees of freedom
AIC: NA

Number of Fisher Scoring iterations: 5
\end{verbatim}
}
\end{frame}

\begin{frame}{Interprétation des paramètres}

Rappelons le modèle :
\begin{eqnarray*}
\lefteqn{\log(E(\texttt{Incidents})) = \hspace*{4cm}}\\
& = & \beta_0 + \log(\texttt{Aggregate}) + \beta_1 \texttt{ship} + \beta_2
\texttt{year}  + \beta_3 \texttt{period}.
\end{eqnarray*}

\bigskip
En fait, chaque facteur correspond à plusieurs paramètres (nombre de
niveaux $-1$). Par exemple, si on considère que les navires de type A sont la
référence, $\beta_1 \texttt{ship}$ est :
$$ \beta_{1B} \iota(\texttt{shipB}) +\beta_{1C} \iota(\texttt{shipC})
+ \beta_{1D} \iota(\texttt{shipD}) + \beta_{1E} \iota(\texttt{shipE}),$$
où $\iota(\texttt{ship}i)$ est une variable indicatrice prenant la valeur
1 pour les navires de type $i$ et 0 sinon.
\end{frame}

\begin{frame}{}
Par conséquent, $\beta_{1B}$ est l'effet sur $\log(\lambda)$ pour les navires de
type B par rapport aux navires de type A. La valeur estimée de
$\log(\lambda)$ pour les navires de type B est inférieure de 0.54334 à celle des navires
de type A (référence).  De manière équivalente, on peut dire que le
nombre espéré estimé de dommages $\lambda$ est multiplié par
$\exp(-0.54334)=0.58$ pour les navires de type B par rapport aux navires de type A,
et est donc réduit.

\bigskip
Les navires de type B et C ont le risque le plus faible, ceux de type E le
plus élevé.

\bigskip
Les navires les plus anciens semblent être les plus sûrs, ceux construits entre
1970-74 ayant le risque le plus élevé.
\end{frame}

\begin{frame}{}
\textbf{Remarque sur les covariables de type facteur} :

Dans l'approche par variables indicatrices utilisée
ci-dessus, une référence doit être définie (parfois de manière très arbitraire). L'effet de ce niveau de référence est fondu dans
l'ordonnée à l'origine. L'interprétation se fait par rapport à la référence. (Ici la
contrainte est $\beta_{1A}=0$).

Par défaut dans R, la première catégorie (ordre alphabétique) est prise comme
référence. Cela peut être changé en redéfinissant le facteur avec la fonction
\texttt{relevel}.

\bigskip
Alternativement, on pourrait utiliser un ensemble différent de
{\color{blue}contrastes}, par exemple ceux où l'effet de chaque
variable est contrasté par rapport à une moyenne globale. L'interprétation des
coefficients dans ce cas est différente. (Ici la contrainte est
$\sum_{i=1}^5 \beta_{1i}=0$).
\end{frame}

\begin{frame}{Approche alternative pour traiter la surdispersion}

Utiliser un modèle différent,
par exemple la loi {\color{blue}binomiale négative} (\textit{negative binomial}) (pour
$y=0,1,2,\dots$) :
$$P(Y_i=y_i) = \frac {\Gamma(y_i+\vartheta)}{\Gamma(\vartheta)\Gamma(y_i+1)}\left(\frac{\vartheta}{\mu_i+\vartheta}\right)^\vartheta
\left(\frac{\mu_i}{\mu_i+\vartheta}\right)^{y_i},$$ 
où $\Gamma(t)= \int_0^{\infty} e^{-x} x^{t-1}dx$ (pour un entier $x$,
$\Gamma(x) = (x-1)!$).

Sous ce modèle $E(Y_i)= \mu_i$ et $Var(Y_i)= \mu_i(\mu_i+1/\vartheta)$.
\end{frame}

\begin{frame}{}

La surdispersion peut aussi survenir avec des données binomiales. On peut soit utiliser un
modèle de type « quasi » avec $Var(Y_i) = \sigma^2 v(\mu_i)$ ou la
loi {\color{blue}bêta-binomiale} (\textit{beta-binomial}) ($y=0, 1, \ldots, m_i$) :
$$P(Y_i=y_i) = \left(
\begin{array}{cc}
m_i\\
y_i
\end{array} \right) \frac{B(\alpha+y_i,m_i+\beta-y_i)}{B(\alpha,\beta)},$$
où $B(u,v) = \int_0^1 x^{u-1}(1-x)^{v-1}dx$.

\bigskip
Ici $E(Y_i) = m_i p_i$ et $Var(Y_i) = m_i p_i (1-p_i) (1 + (m_i-1) \tau^2)$,
où $p_i$ et $\tau$ dépendent de $\alpha$ et $\beta$.

\bigskip
Les fonctions \texttt{glm.nb} (paquet \texttt{MASS}) et
\texttt{negbin,betabin} (paquet \texttt{aod}) permettent d'ajuster ces modèles.

\end{frame}

\subsection{Ce qui peut mal tourner... Séparation (quasi-)complète}
\begin{frame}{Séparation (quasi-)complète en régression binaire}

Si les réponses 0 et 1 peuvent être (presque) séparées sur la base des
valeurs de leurs covariables ({\color{blue}pas de chevauchement, ou séparation quasi-complète} (\textit{quasi-separation})), alors il n'y a pas de solution aux
équations de score pour la régression binaire (ou la solution est
$\beta=\infty$).

\bigskip
\textbf{Exemple} : Analyse des excursions récréatives pour évaluer la sensibilité de
l'utilisation des ressources aux frais d'entrée et aux coûts (\citeNP{Gurm:Triv:1996}).
\end{frame}

\begin{frame}{}
Variables disponibles :
\begin{itemize}
\item TRIPS : Nombre d'excursions en bateau au lac Somerville, Est du Texas, en 1980
\item SO : Classement subjectif de la qualité de l'installation
\item SKI : Goût du répondant pour le ski nautique 
\item I : Revenu -- variable catégorielle 
\item FC3 : Variable indicatrice de coût ; 1 si un droit d'utilisation annuel est payé au lac Somerville 
\item C1 : Coût de déplacement au lac Conroe 
\item C3 : Coût de déplacement au lac Somerville
\item C4 : Coût de déplacement au lac Houston 
\end{itemize}
\end{frame}

\begin{frame}[fragile]{}
{\scriptsize
\begin{verbatim}
> summary(glm((TRIPS>0)~SO+SKI+FC3+I+C1+C3+C4,family=binomial,
data=recreat))

Coefficients:
              Estimate Std. Error z value Pr(>|z|)    
(Intercept)  -2.881479   0.435663  -6.614 3.74e-11 ***
SO            1.437079   0.105799  13.583  < 2e-16 ***
SKI           0.405517   0.329754   1.230  0.21879    
FC3          15.660215 907.987771   0.017  0.98624    
I             0.026479   0.084000   0.315  0.75259    
C1            0.008255   0.027099   0.305  0.76064    
C3           -0.080926   0.017134  -4.723 2.32e-06 ***
C4            0.071124   0.021697   3.278  0.00105 ** 
---
Signif. codes:  0 '***' 0.001 '**' 0.01 '*' 0.05 '.' 0.1 ' ' 1

(Dispersion parameter for binomial family taken to be 1)

    Null deviance: 866.53  on 658  degrees of freedom
Residual deviance: 299.49  on 651  degrees of freedom
AIC: 315.49

Number of Fisher Scoring iterations: 16
\end{verbatim}
}
\end{frame}

\begin{frame}[fragile]{Pourquoi ?}
{\scriptsize
\begin{verbatim}
> with(recreat, table(TRIPS > 0, FC3))
       FC3
          0   1
  FALSE 417   0
  TRUE  229  13
\end{verbatim}
}
\begin{center}
\includegraphics[width = 5.5cm]{figures/recreat.pdf}  
\end{center}

\end{frame}


\begin{frame}{Corrections possibles :}
\begin{itemize}
\item régression logistique exacte \cite{Meht:Pate:1995}.
\item approche de régression robuste \cite{Rous:Chri:2002}.
\item méthode de réduction du biais (\textit{bias reduction}) \cite{Firt:1993,Kosm:Firt:2009} pour
réduire le biais en échantillon fini dans la famille des MLG. Non affectée par le manque de chevauchement.

\bigskip

Une fonction de score corrigée est définie $U^*(\beta) = U(\beta) -\mathcal{J}(\beta) \frac{b_1(\beta)}{n}  $
où $\mathcal{J}(\beta)$ est l'information de Fisher et $b_1(\beta)/n$ est
le premier terme dans l'expression asymptotique du biais de l'estimateur du MV. Cela correspond à une log-vraisem\-blance pénalisée (spécification
canonique)
$$l^*(\beta;y) = l(\beta;y) + \frac{1}{2} \log(|\mathcal{J}(\beta)|).$$

L'utilisation de l'AIC pour la sélection de modèles est
controversée. 
\end{itemize}
\end{frame}

\begin{frame}[fragile]{}
{\scriptsize
\begin{verbatim}
> summary(brglm((TRIPS>0)~SO+SKI+FC3+I+C1+C3+C4,family=binomial,
data=recreat))

Coefficients:
            Estimate Std. Error z value Pr(>|z|)    
(Intercept) -2.76325    0.41610  -6.641 3.12e-11 ***
SO           1.39492    0.10187  13.693  < 2e-16 ***
SKI          0.40403    0.32265   1.252  0.21048    
FC3          1.79847    1.71410   1.049  0.29408    
I            0.02348    0.08187   0.287  0.77424    
C1           0.01878    0.02180   0.861  0.38910    
C3          -0.08293    0.01594  -5.202 1.97e-07 ***
C4           0.06244    0.01877   3.326  0.00088 ***
---
Signif. codes:  0 '***' 0.001 '**' 0.01 '*' 0.05 '.' 0.1 ' ' 1

(Dispersion parameter for binomial family taken to be 1)

    Null deviance: 836.18  on 658  degrees of freedom
Residual deviance: 300.67  on 651  degrees of freedom
Penalized deviance: 258.0502 
AIC:  316.67
\end{verbatim}
}
\end{frame}

\subsection{Ce qui peut mal tourner... Phénomène de Hauck-Donner}
\begin{frame}{Phénomène de Hauck-Donner en régression logistique}

La \texttt{z/t-value} est une approximation de Wald
de la log-vraisemblance pour tester $H_0: \beta_i=0$ et est
parfois trompeuse avec
les MLG binomiaux. En effet, une petite valeur de la \texttt{z/t-value} peut
soit correspondre à une petite statistique du rapport de vraisemblance, soit à une
situation où $|\hat{\beta}_i|$ est grand, l'approximation de Wald
est mauvaise et la statistique du rapport de vraisemblance est grande.
Ces problèmes peuvent survenir dans les cas où les probabilités ajustées sont
extrêmement proches de 0 ou 1.

\bigskip
\bigskip
{\color{red} $\bigstar$ \citeN{Hauc:Donn:1977}}


\end{frame}

\begin{frame}{Remerciements}
Ces diapositives ont été inspirées par le cours GLAM de Eva Cantoni de l'Université de Genève.
\end{frame}

\begin{frame}[allowframebreaks]{Références du chapitre}
\tiny
\printbibliography[heading=none,section=\therefsection]
\end{frame}

\end{refsection}

\end{document}
